% !TeX root = ../Cosmology.tex

%% --------------------------------------------------------
\chapter{Фонова космологія: від метрики FLRW до $\Lambda$CDM}
\label{sec:friedmann}
%% --------------------------------------------------------

Космологія описує Всесвіт як єдиний фізичний об'єкт. Її фундаментальним
відправним пунктом є \textbf{космологічний принцип} (принцип однорідності та ізотропії):
на достатньо великих просторових масштабах ($\gtrsim 100$~Мпк) розподіл
матерії та геометрія простору-часу є трансляційно та ротаційно інваріантними, тобто
не виділяють жодного напрямку і жодної точки. Цей симетрійний постулат
жорстко фіксує єдиний допустимий клас геометрій~--- метрику
Фрідмана--Леметра--Робертсона--Вокера (FLRW)~--- і зводить систему тензорних
рівнянь Айнштайна до двох звичайних диференціальних рівнянь для однієї скалярної
функції часу~--- \textbf{масштабного фактора} $a(t)$.
Фізичний зміст $a(t)$ простий: якщо сьогодні ($t = t_0$) дві галактики
розділені відстанню $d_0$, то в момент $t$ відстань між ними була
$d(t) = a(t)\,d_0/a(t_0)$.
Нормуємо $a(t_0) \equiv 1$: тоді $a < 1$ означає минуле
(Всесвіт менший), $a > 1$~--- майбутнє.

\medskip\noindent%
\textbf{Логічна карта розділу}
% ------------------------------
% Логічна карта розділу 1 (фонова космологія)
% ------------------------------
\begin{center}
\begin{tikzpicture}[roadmap]

\node[box] (FLRW) {Метрика FLRW};
\node[box, right=of FLRW] (Friedmann) {Рівняння\\ Фрідмана};
\node[box, right=of Friedmann] (EOS) {Рівняння\\ стану};
\node[box, right=of EOS] (LCDM) {$\Lambda$CDM};
\node[box, right=of LCDM, fill=red!5, draw=red!70!black] (problems) {Проблеми};

\draw[arrow] (FLRW) -- (Friedmann);
\draw[arrow] (Friedmann) -- (EOS);
\draw[arrow] (EOS) -- (LCDM);
\draw[arrow] (LCDM) -- (problems);

% Підписи зверху
\node[font=\tiny\bfseries, text=blue!70!black, above=0.1cm of FLRW, align=center] {Однорідність\\ та ізотропія};
\node[font=\tiny\bfseries, text=blue!70!black, above=0.1cm of Friedmann, align=center] {Динаміка $a(t)$};
\node[font=\tiny\bfseries, text=blue!70!black, above=0.1cm of EOS, align=center] {$w = p/(\rho c^2)$};
\node[font=\tiny\bfseries, text=blue!70!black, above=0.1cm of LCDM, align=center] {Плоский\\ Всесвіт};
\node[font=\tiny\bfseries, text=red!70!black, above=0.1cm of problems, align=center] {Що не так?};

% Пояснення знизу
\node[font=\scriptsize, gray, below=0.1cm of FLRW] {$ds^2 = -c^2dt^2 + a^2\delta_{ij}dx^idx^j$};
\node[font=\scriptsize, gray, below=0.1cm of Friedmann] {$H^2 = \frac{8\pi G}{3}\rho$};
\node[font=\scriptsize, gray, below=0.1cm of EOS] {$\rho \propto a^{-3(1+w)}$};
\node[font=\scriptsize, gray, below=0.1cm of LCDM] {$\Omega_m + \Omega_\Lambda = 1$};
\node[font=\scriptsize, gray, below=0.1cm of problems, align=center] {плоскість, \\горизонт, збурення};

\end{tikzpicture}
\end{center}

Метрика FLRW у \emph{космологічному (власному) часі}~$t$ для загального випадку
тривимірного простору з постійною кривизною $k = \pm 1, 0$ записується як:
\begin{equation}\label{eq:FLRW_general}
  ds^2 = -c^2 dt^2 + a(t)^2 \left[ \frac{dr^2}{1 - kr^2} + r^2(d\theta^2 + \sin^2\theta\,d\phi^2) \right].
\end{equation}
Для локально плаского Всесвіту ($k = 0$) метрика~\eqref{eq:FLRW_general} набуває спрощеного вигляду:
\begin{equation}\label{eq:FLRW}
  ds^2 = -c^2 dt^2 + a(t)^2\,\delta_{ij}\,dx^i dx^j.
\end{equation}

Підставляючи метричний тензор у рівняння Айнштайна
$G_{\mu\nu} + \Lambda g_{\mu\nu} = \dfrac{8\pi G}{c^4}\,T_{\mu\nu}$ та моделюючи космічний
субстрат як \textbf{ідеальну гідродинамічну рідину}, ми задаємо тензор
енергії-імпульсу у \textbf{супутній системі відліку}%
\footnote{Супутня система відліку~--- це система спостерігача, який
  «пливе» разом із космологічним розширенням, тобто має нульову пекулярну
  швидкість. У такій системі речовина локально перебуває у спокої, і
  тензор $T^{\mu}{}_\nu$ набуває діагонального вигляду.}
у діагональній формі:
\begin{equation*}
  T^\mu{}_\nu = \mathrm{diag}(-\rho c^2,\,p,\,p,\,p).
\end{equation*}
Знак мінус перед $\rho c^2$ є наслідком сигнатури метрики $(-,+,+,+)$:
часова компонента $T^0{}_0 = -\rho c^2$ відповідає густині енергії,
просторові $T^i{}_j = p\,\delta^i{}_j$ --- ізотропному тиску.
Обчислення символів Крістофеля для метрики FLRW та згортка з тензором
Рімана дають ненульові компоненти тензора Айнштайна:
\begin{equation*}
  G^0{}_0 = -\frac{3H^2}{c^2} - \frac{3k}{a^2},
  \qquad
  G^i{}_j = -\frac{1}{c^2}\!\left(2\dot H + 3H^2 + \frac{kc^2}{a^2}\right)\delta^i{}_j,
  \qquad
  H \equiv \frac{\dot a}{a}.
\end{equation*}
Величина $H(t)$ називається \textbf{параметром Хаббла}, вона
є ключовою характеристикою космологічної кінематики і визначає
миттєвий темп розширення Всесвіту.

Підставляючи ці вирази у рівняння Айнштайна, отримуємо систему
\textbf{рівнянь Фрідмана}:
\begin{subequations}\label{eq:friedmann_system}
\begin{align}
  H^2 &= \frac{8\pi G}{3}\,\rho - \frac{kc^2}{a^2} + \frac{\Lambda c^2}{3},
  \label{eq:F1}\\[4pt]
  \frac{\ddot a}{a} &= -\frac{4\pi G}{3}\,\left(\rho + \frac{3p}{c^2}\right) + \frac{\Lambda c^2}{3},
  \label{eq:F2}\\[4pt]
  \dot\rho &= -3H\left(\rho + \frac{p}{c^2}\right).
  \label{eq:conservation_law}
\end{align}
\end{subequations}
Перше~\eqref{eq:F1} --- енергетичний баланс ($(00)$-компонента);
друге~\eqref{eq:F2} --- динаміка прискорення ($(ii)$-компонента).
Третє~\eqref{eq:conservation_law} --- космологічний аналог першого
закону термодинаміки: для однорідної сферичної області об'єму
$V \propto a^3$ воно відтворює $dE = -p\,dV$, де $E = \rho c^2 V$
(детальніше~--- дод.~\ref{app:newton_energy}).
Зауважимо, що~\eqref{eq:conservation_law} не є незалежним: воно
автоматично випливає з~\eqref{eq:F1} і~\eqref{eq:F2} через
тотожність Б'янкі $\nabla_\mu G^{\mu\nu} = 0$.
Система~\eqref{eq:friedmann_system} має рівно \emph{два} незалежних
ступені свободи, але третє рівняння зручніше для безпосереднього
інтегрування і широко використовується надалі.

Для точного математичного моделювання реального Всесвіту ми маємо враховувати одночасний внесок усіх його фізичних компонентів. Прирівнюючи перше рівняння Фрідмана~\eqref{eq:F1} до нуля при $k = 0$ та $\Lambda = 0$, ми визначаємо \textbf{критичну густину}~--- густину, якій відповідає геометрично плоский Всесвіт,~--- та пов'язані з нею безрозмірні параметри густини~$\Omega_i$:
\begin{equation}\label{eq:rho_crit_new}
  \rho_\text{cr}(t) \equiv \frac{3H^2(t)}{8\pi G}, \qquad
  \Omega_i(t) \equiv \frac{\rho_i(t)}{\rho_{\rm cr}(t)}, \qquad
  \Omega_\Lambda(t) \equiv \frac{\Lambda c^2}{3H^2(t)}.
\end{equation}

Поділивши перше рівняння Фрідмана на $H^2(t)$, ми можемо записати його у безрозмірній формі через введені параметри $\Omega_i$. При цьому геометричний член кривизни $-kc^2/(a^2 H^2)$ космологи часто формально трактують як внесок від специфічного «квазіматеріального» флюїду, вводячи \textbf{параметр густини кривизни} $\Omega_k$:
\begin{equation}\label{eq:Omega_k_def}
  \Omega_k(t) \equiv -\frac{kc^2}{a^2(t)H^2(t)}.
\end{equation}
З погляду загальної теорії відносності кривизна $k$ є чистою властивістю просторової геометрії і не є матеріальним джерелом. Проте такий математичний трюк дозволяє звести рівняння Фрідмана для довільної топології простору до елегантного алгебраїчного інваріанта:
\begin{equation}\label{eq:omega_sum}
  \sum \Omega_i(t) + \Omega_k(t) = 1,
\end{equation}
де сума береться за всіма реальними фізичними компонентами (радіація, баріони, холодна темна матерія). Якщо Всесвіт є суворо пласким ($k=0$), то $\Omega_k \equiv 0$, а сума густин речовини та темної енергії точно дорівнює одиниці. Натомість у закритому Всесвіті ($k = +1$, сфера) цей ефективний флюїд має від'ємну густину ($\Omega_k < 0$) та ефективне рівняння стану $w_k = -1/3$, а його густина згасає з розширенням як $\rho_k \propto a^{-2}$.

Використовуючи ці визначення, ми можемо переписати перше рівняння Фрідмана~\eqref{eq:F1} у загальному вигляді критерію енергетичного балансу~\cite{Weinberg2008,Baumann2022}:
\begin{equation}\label{eq:omega_sum_k_integrated}
  \sum_i \Omega_i(t) - 1 = \frac{kc^2}{a^2 H^2} \equiv -\Omega_k(t),
\end{equation}
де $\sum_i \Omega_i(t) = \Omega_m(t) + \Omega_r(t) + \Omega_\Lambda(t)$~---
сума внесків матерії, випромінювання і темної енергії
(космологічна стала трактується як компонента з $w=-1$).
Сучасні дані~\cite{planck_2018} дають $|\Omega_k| < 0.005$,
тобто Всесвіт є геометрично плоским з високою точністю.

\begin{keybox}{Геометрична vs динамічна плоскість}
Важливо розрізняти два питання.
\textbf{Геометричне:} яке значення $k$?
Параметр $k = 0, \pm1$ є глобальною топологічною характеристикою
простору~--- її точне значення невідоме і принципово
не може бути встановлено зі спостережень у межах
скінченного горизонту частинок.
\textbf{Динамічне:} наскільки ця кривина \emph{відчутна}
на доступних нам масштабах?
Це вимірює $\Omega_k(t) = -kc^2/(a^2H^2)$~---
величина, що залежить від часу через $a(t)$.

Якщо $k \neq 0$ але $a^2H^2 \gg c^2$, то $|\Omega_k| \ll 1$~---
кривина існує, але радіус кривини $R_{\rm curv} = a/\sqrt{|k|}$
настільки перевищує горизонт Хаббла $c/H$, що в межах
спостережуваного Всесвіту геометрія є практично евклідовою.
Аналогія: поверхня Землі сферична ($k=+1$), але на масштабі
міста виглядає плоскою.

Тому коли ми кажемо \enquote{Всесвіт плаский}~--- ми маємо на увазі
$|\Omega_k(t_0)| < 0.005$, а не обов'язково $k = 0$.
Різниця принципова концептуально, але не практично:
для всіх розрахунків $k = 0$ є бездоганним
наближенням.
\end{keybox}


%% --------------------------------------------------------
\section{Рівняння стану і закони розширення}
\label{sec:eos}
%% --------------------------------------------------------

Динамічна еволюція кожної космічної компоненти повністю визначається
її \textbf{рівнянням стану} $w \equiv p/(\rho c^2)$. За умови сталого $w$
рівняння збереження~\eqref{eq:conservation_law} інтегрується аналітично:
\begin{equation}\label{eq:rho_evolution}
  \rho(a) = \rho_0 \left(\frac{a_0}{a}\right)^{3(1+w)},
\end{equation}
а підстановка у~\eqref{eq:F1} дає кінематичний закон ($w \neq -1$):
\begin{equation}\label{eq:a_evolution}
  a(t) \propto t^{\,\frac{2}{3(1+w)}}.
\end{equation}

Для повноти картини зведемо основні космологічні епохи до таблиці~\ref{tab:epochs}. У ній наведено лише ті компоненти, які домінують у різні періоди
еволюції Всесвіту. Відсутня, зокрема, \textbf{інфляційна епоха}
($w \approx -1$), оскільки вона передує гарячій стадії і буде
розглянута окремо в наступному розділі~\ref{sec:inflation_kinematics}.

\begin{table}[h!]
\centering\small
\caption{Зведена таблиця космологічних епох: рівняння стану~$w$, закон еволюції густини~$\rho(a)$ та масштабного фактора $a(t)$ часі}
\label{tab:epochs}
\begin{tblr}{
  colspec={Q[l,m]X[l,m]Q[l,m]Q[l,m]Q[l,m]Q[l,m]},
  row{1} = {c, m, font=\bfseries},
}
\toprule
Епоха & Джерело маси-енергії & $w$ & $\rho(a)$ & $a(t)$ \\
\midrule
Радіаційна
  & Фотони, релятивістські нейтрино
  & $1/3$
  & $\propto a^{-4}$
  & $\propto t^{1/2}$ \\
Матеріальна
  & Холодна темна матерія + баріони
  & $0$
  & $\propto a^{-3}$
  & $\propto t^{2/3}$ \\
Сьогодні
  & Темна енергія ($\Lambda$-член)
  & $-1$
  & $\mathrm{const}$
  & $\propto e^{Ht}$ \\
\bottomrule
\end{tblr}
\end{table}

Член кривини $\Omega_k$ формально відповідає ефективному флюїду з $w_k = -1/3$ і густиною $\rho_k \propto a^{-2}$, однак він не є самостійною матеріальною компонентою, а виникає з геометрії простору.

%% --------------------------------------------------------
\section{Конформний час}
\label{sec:horizon}
%% --------------------------------------------------------

Динамічна природа Всесвіту FLRW створює серйозний концептуальний виклик для опису його еволюції: оскільки тканина простору неперервно розширюється, класичний фізичний (космічний) час $t$, який вимірюється годинником локально нерухомого спостерігача, виявляється вкрай незручним для аналізу глобальних процесів. Головна проблема полягає в тому, що швидкість світла $c$ є константою лише локально, тоді як на великих масштабах фотони змушені подорожувати простором, який постійно \enquote{втікає} з-під них. Як наслідок, на стандартних просторово-часових діаграмах $(t, r)$ траєкторії світлових променів (світлові конуси) стають складними викривленими лініями, залежними від еволюції густини енергії. Це суттєво заплутує аналіз причинно-наслідкової структури та причинних горизонтів.

Щоб повернути простору-часу геометричну простоту та інтуїтивну наочність, у космології здійснюють перехід до \emph{конформного часу} $\eta$. Математично конформний час визначається диференціальним співвідношенням, яке повністю компенсує часову зміну масштабного фактора $a(t)$:
\begin{equation}\label{eq:conformal_time_def}
  d\eta = \frac{dt}{a(t)}.
\end{equation}
У цих нових змінних просторово-пласка ($k=0$)\footnote{Вибір $k=0$ обґрунтовується спостереженнями: дані
\emph{Planck}~\cite{planck_2018} свідчать, що геометрія Всесвіту
є плоскою з точністю до $\sim 1\%$.} метрика FLRW набуває елегантного конформно-пласького вигляду:
\begin{equation}\label{eq:conformal_metric}
  ds^2 = a^2(\eta)\bigl(-c^2 d\eta^2 + \delta_{ij}\,dx^i dx^j\bigr).
\end{equation}

Фізично цей новий «космічний годинник» можна уявити так: оскільки $d\eta = dt/a$, в ранню епоху малого $a$ один «тік» конформного годинника охоплює дуже короткий фізичний інтервал $dt = a\,d\eta \ll d\eta$. Конформний час таким чином «розтягує» ранню еволюцію Всесвіту і дозволяє однаково детально описувати як стрімкі процеси ранньої плазми, так і повільну пізню динаміку~--- без зміни форми рівнянь.

\noindent\textbf{Фізичні та математичні переваги конформного формалізму.}
Рівняння~\eqref{eq:conformal_metric} показує, що вся динаміка розширення Всесвіту тепер повністю \enquote{захована} у загальному конформному множнику $a^2(\eta)$, який стоїть перед дужками, тоді як геометрія всередині дужок залишається абсолютно тотожною статичному простору Мінковського. Це дає дві ключові переваги для моделювання космології:
\begin{enumerate}
    \item \textbf{Стабільна причинна структура:} Оскільки для фотонів інтервал руху завжди $ds^2 \equiv 0$, з виразу всередині дужок~\eqref{eq:conformal_metric} прямо випливає закон радіального поширення світла: $c\,d\eta = d\chi$ (де $\chi$~--- супутня просторова координата\footnote{Точніше, $\chi$ --- це радіальна супутня координата, пов'язана з декартовими супутніми координатами $x^i$ співвідношенням $\chi = \sqrt{(x^1)^2 + (x^2)^2 + (x^3)^2}$. У випадку радіального руху світла ($d\theta = d\phi = 0$) просторова частина метрики зводиться саме до $d\chi^2$.}). Це означає, що на просторово-часових діаграмах у координатах $(\eta, \chi)$ світлові конуси завжди залишаються ідеальними прямими лініями під кутом $45^\circ$ (за умови однакового масштабу обох осей, тобто при відкладанні $c\eta$ і $\chi$ в одних одиницях), значно спрощуючи розрахунок причинних горизонтів.
    \item \textbf{Незмінна координатна сітка:} У конформному формалізмі розширення простору відображається як роздування самої координатної сітки. Через це «адреси» (супутні координати $\chi$) об'єктів, які беруть участь у загальному хаббловському потоці, залишаються константними. Реальні фізичні швидкості окремих об'єктів відносно цього фону~--- \emph{пекулярні швидкості} (від англ.\ \textit{peculiar}~--- особливий, власний)~--- відокремлюються від загального розширення і не входять у хаббловський потік. На космологічних масштабах ($\gtrsim 100$~Мпк) пекулярними швидкостями зазвичай нехтують; на локальних~--- вони принципові\footnote{Класичний приклад: галактика Андромеда (M31) наближається до Чумацького Шляху з пекулярною швидкістю ${\sim}110\,\text{км/с}$ завдяки локальному гравітаційному притяганню~--- попри загальне хаббловське розширення.}.
\end{enumerate}

%Розв'язуючи рівняння~\eqref{eq:calH_prime} для сталого $w$, ми знаходимо точні закони еволюції $a(\eta)$ для кожної космологічної епохи (табл.~\ref{tab:epochs}).


%% --------------------------------------------------------
\section{Космологічний червоний зсув та космологічні відстані}
\label{sec:redshift_distances}
%% --------------------------------------------------------

Маючи інструмент конформного часу і чітке розуміння причинної структури,
ми тепер можемо систематично описати як вимірюються відстані
у розширюваному Всесвіті~--- і чому у космології їх існує кілька.

Усі макроскопічні спостереження в астрофізиці пов'язані з динамікою
масштабного фактора $a(t)$ через безрозмірну величину~---
\textbf{космологічний червоний зсув}~$z$. Фотон, випромінений
джерелом у момент $t_e$, приходить до спостерігача в момент $t_0$
з розтягненням довжини хвилі:
\begin{equation}\label{eq:redshift}
  1 + z \equiv \frac{\lambda_{\rm obs}}{\lambda_{\rm em}}
             = \frac{a(t_0)}{a(t_e)}.
\end{equation}
При стандартному нормуванні $a(t_0) \equiv 1$ отримуємо
фундаментальний зв'я\-зок між геометрією та спостереженнями:
\begin{equation}\label{eq:a_z}
  a(t_e) = \frac{1}{1+z}.
\end{equation}
Це рівняння є ключем до космологічних вимірювань: воно дозволяє
замінити теоретичну координату $t_e$, яку неможливо виміряти без
апріорного знання моделі, на безпосередньо спостережуваний
спектроскопічний параметр~$z$.

Реперні значення $z$ в історії Всесвіту:
\begin{align*}
  z &= 0      && \text{(сьогодні, спостерігач)}, \\
  z &\approx 0.3  && \text{(початок прискореного розширення під дією $\Lambda$)}, \\
  z &\approx 1100 && \text{(рекомбінація, формування CMB)}, \\
  z &\sim 10^{10} && \text{(первинний нуклеосинтез)}.
\end{align*}

%% --------------------------------------------------------
\subsection*{Закон Хаббла}
%% --------------------------------------------------------

З відкриття Едвіна Габбла в 1929 році почалася сучасна космологія:
спостерігаючи далекі галактики, він виявив, що їхні спектри зміщені
в червоний бік, причому швидкість віддалення лінійно зростає з відстанню.
Сьогодні ми інтерпретуємо це як наслідок розширення самого простору.

У конформній метриці для плоского Всесвіту~\eqref{eq:conformal_metric}
супутні координати $x^i$ нерухомого спостерігача не змінюються з часом.
Відстань між двома точками вздовж просторової геодезичної в момент $t$
дорівнює добутку масштабного фактора на різницю їхніх супутніх координат:
$d_p(t) = a(t)\,\Delta \chi$. Позначаючи $d_C \equiv \Delta \chi$ (супутня
відстань), маємо:
\begin{equation}
d_p(t) = a(t)\,d_C.
\end{equation}

Диференціюючи цей вираз за часом, отримуємо:
\begin{equation}\label{eq:hubble_law}
\dot{d}_p = \frac{\dot{a}}{a}\,d_p = H(t)\,d_p, \qquad
H \equiv \frac{\dot{a}}{a}.
\end{equation}
Це і є \textbf{закон Хаббла} в диференціальній формі: швидкість
віддалення об'єкта пропорційна його відстані, а коефіцієнтом
пропорційності виступає параметр Хаббла.

Сьогодні ($t = t_0$) маємо $\dot{d}_p = H_0\,d_p$, де
$H_0 = 100\,h\,\text{км/с/Мпк}$ ($h = 0.6736$ за даними
\textit{Planck}~2018~\cite{planck_2018}).

Фізична відстань $d_p(t)$ є \emph{практично
неспостережуваною}. Вона змінюється щомиті через розширення простору,
тоді як світло фіксує об'єкт у минулому, а не в його теперішньому
положенні. Тому закон Хаббла в такому вигляді є лише кінематичним
співвідношенням, а не робочою формулою для вимірювань.

%% --------------------------------------------------------
\subsection*{Супутня відстань}
%% --------------------------------------------------------

Для фотона $ds^2 = 0$, тому з конформної метрики~\eqref{eq:conformal_metric}
випливає $c\,d\eta = d\chi$ (див.\ §\,\ref{sec:horizon}). Отже, супутня
відстань до об'єкта, що випромінив світло в момент $\eta_e$ (який ми
спостерігаємо в момент $\eta_0$), дорівнює:
\begin{equation*}
d_C \equiv \chi = c\int\limits_{\eta_e}^{\eta_0} d\eta'
           = c\int\limits_{t_e}^{t_0} \frac{dt}{a(t)}.
\end{equation*}

Величина $c\,dt/a(t)$ має просту інтерпретацію: це елемент
\textbf{супутньої відстані}, який світло проходить за фізичний час $dt$.
Обернена величина $a(t)/c$ визначає миттєвий \textbf{горизонт Хаббла}
$R_H(t) = c/H(t)$ (у фізичних координатах), який відокремлює області, що віддаляються повільніше або швидше
за швидкість світла.


\begin{keybox}{Конформний горизонт Хаббла}
У конформній метриці \eqref{eq:conformal_metric} природно ввести
конформний параметр Хаббла $\mathcal{H} \equiv a'/a = aH$ (штрих ---
похідна за $\eta$). Тоді конформний горизонт Хаббла дорівнює
\begin{equation*}
    c\mathcal{H}^{-1} = \frac{c}{aH}.
\end{equation*}

На відміну від звичайного радіуса Хаббла $R_H = c/H$, конформний горизонт $c\mathcal{H}^{-1}$ вимірюється в супутніх координатах. Це робить його зручним інструментом для порівняння фізичних масштабів космологічних структур (наприклад, BAO) із причинною структурою Всесвіту в будь-який момент часу.
\end{keybox}


Оскільки безпосередньо виміряти космічний час $t_e$ неможливо,
перейдемо до спостережуваного червоного зсуву $z$. Диференціюючи~\eqref{eq:a_z} отримуємо:
\begin{equation*}
\dot{a} = -\frac{1}{(1+z)^2}\,\dot{z}.
\end{equation*}

Підставляючи це у визначення параметра Хаббла $H \equiv \dot{a}/a$, маємо:
\begin{equation*}
    H = \frac{\dot{a}}{a} = (1+z) \cdot \left( -\frac{\dot{z}}{(1+z)^2} \right)
    = -\frac{1}{1+z}\,\frac{dz}{dt},
\end{equation*}
звідки знаходимо потрібний диференціал:
\begin{equation*}
d\eta = \frac{dt}{a(t)} = -\frac{dz}{H(z)}.
\end{equation*}

Далі необхідно знати явну залежність $H(z)$. Для цього розділимо
перше рівняння Фрідмана~\eqref{eq:F1} на $H_0^2$. Використовуючи безрозмірні параметри густини
$\Omega_i \equiv \rho_{i0}/\rho_{\rm cr0}$ та закон еволюції кожної
компоненти $\rho(a) = \rho_0 (a_0/a)^{3(1+w)}$ (див.~\eqref{eq:rho_evolution}),
отримуємо:
\begin{equation*}
    \frac{H^2(z)}{H_0^2} = \Omega_r(1+z)^4 + \Omega_m(1+z)^3 + \Omega_k(1+z)^2 +  \Omega_\Lambda,
\end{equation*}
де $a_0/a = 1+z$. Таким чином, безрозмірна функція Хаббла
$E(z) \equiv H(z)/H_0$ має вигляд:
\begin{equation}\label{eq:Ez}
E(z) = \sqrt{\Omega_r(1+z)^4 + \Omega_m(1+z)^3 + \Omega_k(1+z)^2 + \Omega_\Lambda}.
\end{equation}

При переході від $t$ до $z$ межі інтегрування змінюються: моменту
випромінювання $t_e$ відповідає червоний зсув $z$, а сучасному
моменту $t_0$ --- $z=0$. Зміна знаку $dz$ компенсує переворот меж
($\int_z^0 \to -\int_0^z$). Остаточно отримуємо:
\begin{equation}\label{eq:comoving_dist}
d_C(z) = \frac{c}{H_0}\int\limits_0^z \frac{dz'}{E(z')}.
\end{equation}

Оскільки аналітично цей інтеграл для комбінації різних степенів
$(1+z)$ під коренем взяти не можна, після знаходження оптимальних
параметрів $\Omega_i$ розрахунок реальних відстаней за формулою
\eqref{eq:comoving_dist} виконується виключно чисельними методами.

%% --------------------------------------------------------
\subsection*{Система космологічних відстаней}
%% --------------------------------------------------------

У статичному просторі відстань між двома точками однозначна.
У розширюваному Всесвіті FLRW це не так: кожна фізична задача
(вимірювання потоку, кутового розміру, швидкості розлітання)
потребує своєї «лінійки». Спостерігач може безпосередньо виміряти
лише три величини~--- червоний зсув $z$, потік енергії $F$ та
кутовий розмір $\delta\theta$,~--- тому вводять систему
\textbf{оперативних відстаней}, кожна з яких відповідає одному
типу вимірювань.

%% --------------------------------------------------------
\subsection*{Відстань люмінозності~--- стандартні свічки}
%% --------------------------------------------------------

Для об'єкта відомої світності $L$ у статичному просторі:
$F = L/(4\pi d^2)$. У розширюваному Всесвіті потік послаблюється
в $(1+z)^2$ разів через два ефекти:
\begin{enumerate}
  \item кожен фотон втрачає енергію: $E_{\rm obs} = E_{\rm em}/(1+z)$;
  \item темп прибуття фотонів сповільнюється: $\Delta t_{\rm obs}
        = \Delta t_{\rm em}(1+z)$.
\end{enumerate}
Щоб зберегти форму $F = L/(4\pi d_L^2)$, вводять:
\begin{equation}\label{eq:d_L}
  d_L(z) = d_C(z)\,(1+z).
\end{equation}
Саме вимірювання $d_L(z)$ для наднових типу~Ia у 1998~р.\ призвело
до відкриття прискореного розширення та темної
енергії~\cite{Perlmutter1999,Riess1998}.

%% --------------------------------------------------------
\subsection*{Кутова відстань~--- стандартні лінійки}
%% --------------------------------------------------------

\textbf{Кутова відстань} визначається через видимий кутовий розмір
$\delta\theta$ об'єкта з відомим фізичним розміром~$\ell$:
\begin{equation}\label{eq:d_A_def}
  d_A \equiv \frac{\ell}{\delta\theta}.
\end{equation}
Оскільки кут фіксується геометрією епохи випромінювання, коли
масштабний фактор був $a_e = 1/(1+z)$, то:
\begin{equation}\label{eq:d_A}
  d_A(z) = \frac{d_C(z)}{1+z}.
\end{equation}
Зауважимо, що $d_A(z)$ \emph{не є монотонною}: вона досягає
максимуму при $z \approx 1.6$ (для параметрів $\Lambda$CDM) і
далі спадає~--- об'єкти при дуже великих $z$ виглядають
\emph{більшими}, ніж при проміжних. Фізично це відбувається тому,
що $d_A = d_C/(1+z)$: при малих $z$ чисельник $d_C(z)$ зростає
швидше ніж знаменник $(1+z)$, а при великих $z$ знаменник починає
«перемагати». Конкретне значення $z \approx 1.6$ визначається
балансом $\Omega_m$ і $\Omega_\Lambda$ і є характеристикою
саме нашого Всесвіту.
Цей ефект «космологічного
збільшення» детально розглядається у задачі~\ref{prob:d_A_max}.

На практиці $d_A(z)$ є критичним інструментом для аналізу
кутового масштабу акустичних піків CMB та BAO~--- саме через неї
вимірюють геометрію Всесвіту та $H_0$.

Відстані $d_A(z)$ та $d_L(z)$ жорстко пов'язані між собою через супутню відстань $d_C(z)$ та \textbf{теорему взаємності Етерінгтона}:
\begin{equation}\label{eq:etherington}
  d_L(z) = (1+z)^2\, d_A(z).
\end{equation}
Цє співвідношення є строгим геометричним інваріантом будь-якої метричної теорії гравітації. Воно виконується за умови збереження кількості фотонів у геометричній оптиці та є прямим наслідком теореми Ліувілля про збереження фазового об'єму вздовж світлових променів~\cite{Baumann2022}. На практиці рівність~\eqref{eq:etherington} слугує фундаментальним тестом на відсутність екзотичної фізики (наприклад, осциляцій фотонів в аксіоноподібні частинки) та міжгалактичного поглинання світла, а також дозволяє безпосередньо перераховувати кутові масштаби джерел у їхні фотометричні характеристики.


%% --------------------------------------------------------
\section{Параметри густини та модель $\Lambda$CDM }
\label{sec:lcdm_and_flatness}
%% --------------------------------------------------------

Параметри $\Omega_i$ неможливо виміряти безпосередньо ---
вони визначаються через спільний чисельний фіт спостережних даних
різних типів.

\begin{enumerate}
    \item \textbf{Спостереження:} для кожного об'єкта вимірюють
    червоний зсув $z$ та одну або кілька залежних від космології
    величин:
    \begin{itemize}
        \item для наднових типу Ia --- видиму зоряну величину $m(z)$;
        \item для BAO --- кутовий розмір $\delta\theta(z)$ та/або
        інтервал червоного зсуву $\Delta z(z)$;
        \item для CMB --- спектр потужності $C_\ell$.
    \end{itemize}

    \item \textbf{Вибір моделі:} фіксують набір параметрів
    $\{H_0, \Omega_m, \Omega_\Lambda, \Omega_k, \Omega_b, \dots\}$
    та функцію $E(z) = \sqrt{\Omega_r(1+z)^4 + \Omega_m(1+z)^3
    + \Omega_k(1+z)^2 + \Omega_\Lambda}$.

    \item \textbf{Пряме обчислення теоретичних передбачень:}
    для заданого набору параметрів чисельно обчислюють:
    \begin{equation*}
    d_C(z) = \frac{c}{H_0} \int\limits_0^z \frac{dz'}{E(z')},
    \qquad
    d_L(z) = (1+z)\,d_C(z),
    \qquad
    d_A(z) = \frac{d_C(z)}{1+z},
    \end{equation*}
    а також, за потреби, звуковий горизонт $r_s(z_*)$ та
    теоретичний спектр $C_\ell$.

    \item \textbf{Порівняння з даними:} обчислюють $\chi^2$ як
    суму квадратів відхилень між спостереженими величинами
    ($m(z)$, $\delta\theta(z)$, $C_\ell$, \dots) та їхніми
    теоретичними виразами, наприклад:
    \begin{equation*}
    m^{\text{th}}(z) = M + 5\log_{10}\!\left(\frac{d_L(z)}{10\ \text{пк}}\right),
    \qquad
    \delta\theta^{\text{th}}(z) = \frac{r_s}{d_A(z)}.
    \end{equation*}

    \item \textbf{MCMC-фітинг:} MCMC перебирає мільйони комбінацій
    параметрів, повторюючи кроки 3--4 для кожної, і знаходить ті,
    що мінімізують $\chi^2$ --- це і є найкращі значення
    $\Omega_i$, $H_0$ тощо.
\end{enumerate}

Зокрема, за сучасними прецизійними даними \textit{Planck} 2018~\cite{planck_2018}, результатом такого фітингу для базової пласкої моделі нашого Всесвіту є наступні фундаментальні константи:
\begin{equation}\label{eq:LCDM_params}
\begin{split}
    \Omega_m &= 0.3153 \pm 0.0073,        \\
    \Omega_\Lambda &= 0.6847 \pm 0.0073,  \\
    \Omega_r &\approx 9.4 \times 10^{-5}, \\
    \Omega_k &= 0.0007 \pm 0.0019.
\end{split}
\end{equation}

Ті самі параметри дозволяють обчислити \textbf{вік Всесвіту}. З першого рівняння Фрідмана $dt = da/(aH)$, тому:
\begin{equation}\label{eq:age}
  t_0 = \int\limits_0^1 \frac{da}{a\,H(a)} = \frac{1}{H_0}\int\limits_0^1 \frac{da}{a\,E(a)},
\end{equation}
де $E(a) = \sqrt{\Omega_r a^{-4} + \Omega_m a^{-3} + \Omega_\Lambda}$~--- функція~\eqref{eq:Ez}, переписана через закон зв'язку масштабного фактора і червоного зсуву $a = 1/(1+z)$. Оскільки аналітичного первісного виразу для такої комбінації не існує, інтеграл обчислюється чисельними методами. Підставляючи параметри \textit{Planck} 2018, отримуємо:
\begin{equation*}
  t_0 = 13.801 \pm 0.024 \text{ млрд років.}
\end{equation*}

Тепер ми маємо повне право звести всі компоненти нашого плаского Всесвіту в єдину фінальну балансову таблицю~\ref{tab:lcdm_components}.


\begin{table}[ht]
\centering\small
\caption{Компоненти стандартної моделі $\Lambda$CDM та їхні сучасні параметри густини за даними місії \textit{Planck}~\cite{planck_2018}.}
\label{tab:lcdm_components}
\begin{tblr}{
  colspec={Q[l,m, 2.5cm]Q[c,m]Q[c,m]QX[l, m]},
  row{1}={font=\bfseries, c, m},
}
\toprule
Компонента & Параметр & $w$ & {Значення\\при $z=0$} & {Фізична природа} \\
\midrule
Баріони                   & $\Omega_b$        & $0$   & {$\approx 0.0493$\\ ($\approx 5\%$)}   & Звичайна речовина (протони, електрони, атоми)   \\
Холодна темна матерія     & $\Omega_c$        & $0$   & {$\approx 0.2660$\\ ($\approx 27\%$)}  & Небаріонні слабковзаємодіючі масивні частинки   \\
Темна енергія ($\Lambda$) & $\Omega_\Lambda$  & $-1$  & {$\approx 0.6847$\\ ($\approx 68\%$)}  & Енергія космологічного вакууму                  \\
Випро\-мінюван\-ня            & $\Omega_r$        & $1/3$ & {$\approx 0.000094$\\ ($\approx 0.009\%$)} & Реліктові фотони та космологічні нейтрино       \\
\bottomrule
\end{tblr}
\end{table}



%% --------------------------------------------------------
\section{Парадокс плоскості: навіщо таке точне налаштування?}
\label{subsec:flatness_paradox}
%% --------------------------------------------------------

Модель $\Lambda$CDM вказує на те, що сьогодні $\Omega_\text{tot}(t_0) \approx 1$. Здавалося б, це свідчить про геометричну простоту нашого Всесвіту. Проте, якщо екстраполювати рівняння Фрідмана у зворотному часовому напрямку і проаналізувати, якою мала бути ця величина на ранніх етапах еволюції, результат виявляється концептуально парадоксальним.

Проаналізуємо, як відхилення від плоскості поводилося в минулому.
З рівняння~\eqref{eq:omega_sum_k_integrated} видно, що воно визначається
комбінацією $a^2 H^2$ у знаменнику. Простежимо її долю в кожну епоху:
\begin{itemize}
  \item \textbf{Радіаційно-домінантна епоха ($w = 1/3$):} Оскільки $H^2 \propto \rho_r \propto a^{-4}$, маємо $a^2 H^2 \propto a^{-2}$. Тоді відхилення від плоскості прогресує лінійно по часі:
  \begin{equation}\label{eq:omega_evolution_r_new}
    |\Omega_{\rm tot}(t) - 1| \propto |k|a^2(t) \propto |k|t.
  \end{equation}
  \item \textbf{Матеріально-домінантна епоха ($w = 0$):} Тут $H^2 \propto \rho_m \propto a^{-3}$, відповідно $a^2 H^2 \propto a^{-1}$. Відхилення від плоскості продовжує зростати:
  \begin{equation}\label{eq:omega_evolution_m_new}
    |\Omega_{\rm tot}(t) - 1| \propto |k|a(t) \propto |k|t^{2/3}.
  \end{equation}
\end{itemize}

Обидва рівняння говорять одне: якщо $k \neq 0$ (тобто Всесвіт не є ідеально плоским), то з часом відхилення $|\Omega_{\rm tot}(t)-1|$ не зникає, а невпинно зростає. Плоский Всесвіт є точкою нестійкої рівноваги: будь-яке початкове відхилення від критичної густини призводить до дедалі більшого відхилення в майбутньому.

%Обидва рівняння говорять одне: $\Omega_{\rm tot} = 1$ є точкою
%\emph{нестійкої рівноваги}~--- як олівець, поставлений на гострий кінець.
%Будь-яке початкове відхилення від критичної густини не затухає,
%а невпинно наростає з часом.

Щоб отримати спостережуване сьогодні значення $|\Omega_k(t_0)| < 0.005$
після того, як Всесвіт охолонув від температури епохи
Великого Об'єднання ($T_{\rm GUT} \sim 10^{26}$~К, масштаб енергій
${\sim}10^{16}$~ГеВ) до сучасного стану ($T_0 \approx 2.73$~К)~---
оскільки для релятивістської речовини $a \propto T^{-1}$
(випливає з $\rho_r \propto a^{-4}$ та закону Стефана--Больцмана
$\rho_r \propto T^4$), це відповідає розширенню масштабного фактора
у $a_0/a_{\rm GUT} = T_{\rm GUT}/T_0 \approx 10^{26}$ разів~--- початкові умови густини
на самому старті мали бути налаштовані з фантастичною,
абсолютно неприродною точністю:
\begin{equation*}
  |\Omega_{\rm tot} - 1|_{t = t_{\rm Pl}} \lesssim 10^{-60}.
\end{equation*}
Такий ступінь узгодженості початкових умов свідчить про наявність глибокої
концептуальної проблеми в межах класичної Фрідманівської космології. Відхилення
від критичної густини у планківську епоху навіть на величину порядку $10^{-60}$
призвело б до кардинальної зміни космологічного сценарію: у випадку замкненої
геометрії Всесвіт зазнав би колапсу назад у сингулярність задовго до початку
епохи нуклеосинтезу, а у випадку відкритої~--- надто швидке розширення унеможливило
б гравітаційну нестабільність і подальше формування великомасштабної структури.

Фізики називають цю ситуацію \textbf{проблемою тонкого налаштування}
(\emph{fine-tuning problem}).
Саме по собі число $10^{-60}$ не заборонено жодним законом природи.
Але наука не приймає «так склалося» як відповідь~--- особливо
коли точність потрібна з шістдесятьма знаками після коми.

\begin{keybox}{Динамічне вирішення: космічна інфляція~\cite{linde_1982,albrecht_1982}}

У 1980--1981 роках Алан Гут запитав: а що, якщо до гарячої стадії
Всесвіт пережив коротку фазу \emph{експоненціального} розширення?
Тоді закон зміни $a^2H^2$ був би оберненим.
Під час інфляції домінує скалярне поле з $w \approx -1$,
параметр Хаббла залишається майже сталим,
а масштабний фактор летить вгору як $a(t) \propto e^{Ht}$.

У цей період знаменник $a^2H^2$ стрімко зростає, а відхилення від плоскості зазнає потужного придушення:
\begin{equation}
  |\Omega_{\rm tot}(t) - 1| = \frac{kc^2}{a^2 H^2} \propto \frac{1}{a^2(t)} \propto e^{-2Ht}.
\end{equation}
Саме стільки потрібно, щоб \enquote{згладити} довільне початкове відхилення
$|\Omega-1| \sim 1$ (природне значення на планківській епосі) до рівня,
сумісного з подальшою еволюцією: після $N = 60$ $e$-folds отримуємо
$|\Omega-1| \lesssim e^{-2 \times 60} \approx 10^{-52}$.
Цього з великим запасом достатньо, щоб сьогодні спостерігати
$|\Omega_k| < 0.005$~\cite{planck_2018}~--- адже між кінцем інфляції
і сьогоднішнім днем відхилення від плоскості вже \emph{зростало}
за законом~\eqref{eq:omega_sum_k_integrated}, проте стартувало з
настільки малого значення, що навіть це зростання не встигло його
підняти до помітного рівня.
\end{keybox}


%% --------------------------------------------------------
\section*{Підсумок і відкриті питання}
\addcontentsline{toc}{section}{Підсумок і відкриті питання}
%% --------------------------------------------------------

\textbf{Що ми збудували.}
Починаючи з одного постулату~--- однорідності та ізотропії~--- ми отримали
повний кінематичний і динамічний опис однорідного Всесвіту, що розширюється.
Його серцевина~--- система рівнянь Фрідмана~\eqref{eq:friedmann_system}~---
зводить задачу до єдиної функції $a(t)$, еволюція якої повністю визначається
складом речовини через $\rho(t)$ та $p(t)$.
Модель $\Lambda$CDM~--- шість вільних параметрів, чотири компоненти,
$\Omega_{\rm tot} = 1$~--- є найкращим з існуючих описів спостережуваного
Всесвіту на великих масштабах~\cite{planck_2018}.

\medskip
\textbf{Три проблеми, які $\Lambda$CDM не вирішує.}
Попри видатний емпіричний успіх, ця модель є теорією \emph{еволюції},
а не \emph{початкових умов}.
Вона не відповідає на три фундаментальні питання:

\begin{enumerate}
  \item \textbf{Проблема плоскості.}
  Чому $|\Omega_k(t_0)| < 0.005$, якщо плоска геометрія є точкою
  динамічної нестабільності~\eqref{eq:omega_sum_k_integrated} і будь-яке
  початкове відхилення від неї мало б лавиноподібно зростати?
  За даними Планківської епохи це вимагає тонкого налаштування
  $|\Omega_{\rm tot} - 1|_{t_{\rm Pl}} \lesssim 10^{-60}$.

  \item \textbf{Проблема горизонту.}
  Реліктове випромінювання є однорідним з точністю до $10^{-5}$ по всьому
  небу. Але в \emph{стандартній гарячій космології без інфляції} ділянки,
  розділені більш ніж на ${\sim}2°$, ніколи не перебували у причинному
  контакті: їхні горизонти Хаббла $\mathcal{H}^{-1}$ на момент рекомбінації
  не перетиналися, тобто між ними не міг пройти жодний сигнал від Великого
  вибуху до рекомбінації. Але якщо ці ділянки ніколи не «спілкувалися»~---
  чому вони мають однакову температуру з точністю $10^{-5}$?
  Термодинамічна рівновага без причинного контакту неможлива~---
  це і є проблема горизонту.

  \item \textbf{Проблема первинних збурень.}
        У межах стандартної моделі $\Lambda$CDM наявність початкових збурень
        густини, які стали зародками всієї великомасштабної структури Всесвіту,
        приймається як суто феноменологічна початкова умова. Класична Фрідманівська
        космологія принципово неспроможна пояснити мікрофізичну природу виникнення
        цих неоднорідностей, а також причини формування їхнього \emph{майже строго
        степеневого спектра з надзвичайно малою амплітудою}, залишаючи це питання
        відкритим до залучення механізмів квантової генерації на стадії інфляції%
        \footnote{До цього переліку іноді додають \textbf{проблему ентропії}:
        спостережуваний Всесвіт має колосальну ентропію $S \sim 10^{88}$~---
        звідки вона взялась? Розігрів після інфляції дає природну відповідь:
        вся ентропія народжується при термалізації продуктів розпаду інфлатона.}.
\end{enumerate}

\medskip
Усі три зазначені космологічні парадокси мають спільне елегантне вирішення. Якщо припустити, що гарячій стадії Великого вибуху передувала \emph{епоха інфляції}~--- надшвидкого розширення Всесвіту, керованого особливим скалярним полем,~--- то всі суперечності знімаються одночасно. Геометрична плоскість стає наслідком колосального розтягнення простору; температурна однорідність (проблема горизонту) пояснюється тим, що весь спостережуваний Всесвіт розширився з крихітної причинно-зв'язаної області; а первинні неоднорідності виявляються розтягнутими до космічних масштабів квантовими коливаннями цього самого поля. Детальному математичному опису та аналізу цього дивовижного механізму присвячений наступний розділ.

% ============================================================
% Задачі до розділу 1
\section*{Задачі}
\addcontentsline{toc}{section}{Задачі}
% ============================================================

\begin{problem}{Точна підгонка: закон Габбла в метриці FLRW}
    У законі Габбла $\dot{d}_p = H(t) d_p$ відстань $d_p$ є фізичною
    (власною) відстанню. Нехай спостерігач сьогодні ($t=t_0$) реєструє
    фотон від галактики, який був випромінений у момент $t_e$.
    \begin{enumerate}
        \item Виразіть $d_p(t_e)$ через $d_C$ та $a(t_e)$.
        \item Виразіть $d_p(t_0)$ через $d_C$ та $a(t_0)$.
        \item Покажіть, що відношення $\dot{d}_p(t_0)$ до $d_p(t_0)$
              дає $H_0$, але чому це не можна безпосередньо виміряти
              для далекої галактики?
    \end{enumerate}
\end{problem}

\begin{problem}{Інтегрування $d_C(z)$ для Всесвіту без $\Lambda$}
    Розглянемо плоский Всесвіт, заповнений лише нерелятивістською
    матерією ($\Omega_m = 1$, $\Omega_\Lambda = \Omega_r = 0$).
    \begin{enumerate}
        \item Знайдіть явний вигляд функції $E(z)$.
        \item Обчисліть аналітично $d_C(z)$ за формулою \eqref{eq:comoving_dist}.
        \item Знайдіть граничне значення $d_C(z)$ при $z \to \infty$.
              Чи є воно скінченним? Чому?
    \end{enumerate}
\end{problem}

\begin{problem}{Вік Всесвіту в різних моделях}
    Використовуючи співвідношення $dt = da/(aH)$ та закон еволюції
    $a(t) \propto t^{2/[3(1+w)]}$ для домінуючої компоненти:
    \begin{enumerate}
        \item Виразіть вік Всесвіту $t_0$ через $H_0$ в моделях, де
              домінує випромінювання ($w=1/3$) або матерія ($w=0$).
        \item Для $\Lambda$CDM параметри $\Omega_m \approx 0.315$, $\Omega_\Lambda \approx 0.685$.
              Поясніть, чому $t_0$ більший, ніж у модельному випадку
              з $w=0$ при тому самому $H_0$.
    \end{enumerate}
\end{problem}

\begin{problem}{Проблема плоскості в числах}
    Сьогодні $|\Omega_k(t_0)| < 0.005$.
    У радіаційно-домінантну епоху $|\Omega_k(t)-1| \propto a^2(t) \propto t$.
    \begin{enumerate}
        \item Оцініть $|\Omega_k - 1|$ на момент нуклеосинтезу
              ($t_{\text{BBN}} \approx 1$ с), використовуючи
              $t_0 \approx 13.8$ млрд років.
        \item Оцініть $|\Omega_k - 1|$ на планківському часі
              ($t_{\text{Pl}} \approx 5.4 \times 10^{-44}$ с).
        \item Чому число $10^{-60}$ (згадане в підрозділі~\ref{subsec:flatness_paradox})
              виникає природно? (Підказка: $t_0 / t_{\text{Pl}} \sim 10^{60}$.)
    \end{enumerate}
\end{problem}

\begin{problem}{Точка повороту кутової відстані}\label{prob:d_A_max}
    Дослідіть аналітично ефект максимуму кутової відстані
    («космологічного збільшувального скла»).
    Для спрощення розглянемо Всесвіт Ейнштейна--де-Сіттера
    (плоский, $\Omega_m = 1$, $\Omega_\Lambda = 0$).
    \begin{enumerate}
        \item Покажіть, що $E(z) = (1+z)^{3/2}$.
        \item Знайдіть явний аналітичний вираз для $d_C(z)$ та $d_A(z)$.
        \item Визначте червоний зсув $z_{\text{max}}$, при якому $d_A(z)$
              досягає максимуму.
        \item Обчисліть кутовий розмір $\delta\theta$ галактики довжиною
              $\ell = 10$ кпк, яка перебуває на $z = z_{\text{max}}$ та на $z = 4$.
              Порівняйте їх та зробіть фізичний висновок.
    \end{enumerate}
\end{problem}

\begin{problem}{Чисельне обчислення $d_C(z)$ та фіт параметрів $\Lambda$CDM за даними BAO}
\label{prob:bao_fit_program}

У цій задачі вам пропонується самостійно реалізувати чисельний
фіт космологічних параметрів за спостереженнями BAO.

\textbf{Дані BAO.} У таблиці~\ref{tab:bao_data} наведено
вимірювання величини $D_V(z)/r_s$ для різних червоних зсувів $z$,
де $D_V(z)$ --- усереднена відстань:
\begin{equation*}
    D_V(z) = \left[ z \, d_C^2(z) \, \frac{c}{H(z)} \right]^{1/3},
\end{equation*}
а $r_s = 147.09\ \text{Мпк}$ --- звуковий горизонт (фіксовано за
даними \textit{Planck} 2018).

Спостережні дані BAO: червоний зсув $z$, значення $D_V/r_s$ та похибка $\sigma$
\begin{center}
\label{tab:bao_data}
\begin{tblr}{
  colspec={Q[c,m]Q[c,m]Q[c,m]},
  row{1} = {font=\bfseries},
}
\toprule
$z$ & $D_V/r_s$ & $\sigma$ \\
\midrule
0.106 & 3.047 & 0.137 \\
0.150 & 4.480 & 0.168 \\
0.320 & 8.467 & 0.167 \\
0.570 & 13.773 & 0.134 \\
0.440 & 11.550 & 0.400 \\
0.600 & 14.940 & 0.420 \\
0.730 & 16.850 & 0.600 \\
1.520 & 26.130 & 0.580 \\
\bottomrule
\end{tblr}
\end{center}

\textbf{Завдання.}

\begin{enumerate}
    \item Напишіть програму (мовами Python, Julia, C++ або іншою на ваш вибір),
          яка:
          \begin{itemize}
              \item визначає функцію $E(z)$ та функцію для обчислення $d_C(z)$
                    з використанням чисельного інтегрування (наприклад, метод
                    Сімпсона, Гауса або вбудований інтегратор);
              \item обчислює теоретичне значення $D_V(z)/r_s$ для заданих
                    $H_0$ та $\Omega_m$;
              \item реалізує $\chi^2$ функцію:
                    \[
                        \chi^2(H_0, \Omega_m) = \sum_{i} \left( \frac{ (D_V/r_s)_{\text{th}}(z_i) - (D_V/r_s)_{\text{obs},i} }{\sigma_i} \right)^2;
                    \]
              \item знаходить мінімум $\chi^2$ за допомогою методу
                    сітки, градієнтного спуску або вбудованого оптимізатора
                    (наприклад, \texttt{scipy.optimize.minimize}).
          \end{itemize}

    \item Знайдіть оптимальні значення $H_0$ та $\Omega_m$,
          побудуйте графік $D_V(z)/r_s$ (теорія разом з даними).

    \item Обчисліть вік Всесвіту $t_0$ для отриманих параметрів.

    \item \textbf{Дослідження.} Змініть модель, дозволивши $\Omega_k \neq 0$
          (тобто $\Omega_\Lambda$ і $\Omega_m$ --- вільні, а $\Omega_k = 1 - \Omega_m - \Omega_\Lambda - \Omega_r$).
          Чи покращується $\chi^2$? Чи сумісне отримане $\Omega_k$ з нулем?
\end{enumerate}

\textbf{Вказівки:}
\begin{itemize}
    \item Швидкість світла $c = 299792.458\ \text{км/с}$.
    \item Пошук мінімуму можна почати з початкового наближення
          $H_0 \approx 70\ \text{км/с/Мпк}$, $\Omega_m \approx 0.3$.
    \item Для чисельного інтегрування можна використовувати метод
          Сімпсона або вбудовану функцію \texttt{quad} з \texttt{scipy.integrate}
          (Python), \texttt{quadgk} (Julia) або аналогічну.
\end{itemize}
\end{problem}